\section{Conclusion}
\label{sec:conclusion}

The joint-anchored interaction field of SHOW3D is a deterministic function of the hand joints, the object pose and a known mesh, and this paper asked what is gained by predicting those arguments and applying the function rather than regressing the field directly. Under a protocol fixed in advance, at a frozen representation shared by every model, the factored formulation lowers the held-out error of a matched direct regressor from 46.2 to 44.4\,mm on the official development subjects and from 35.7 to 34.1\,mm on a second pair, with two seeds whose ranges do not overlap, and a single-view factored model outperforms the two-view direct one. The gain grows with the difficulty of the target, from under a millimetre at contact to 21\,mm when the hand leaves the image. Opening the model showed that the gain is carried by the direct head and delivered through the shared representation: the geometric branch is a structured auxiliary task, inference-time fusion contributes a tenth of a millimetre, and the branch can be dropped at deployment. The object pose is the bottleneck of the geometric route in domain and the joints are the bottleneck out of domain, and the advantage does not survive a change of dataset or the addition of a second one to training. We release the code, the derived labels, the HOT3D-IF protocol, every checkpoint and the per-frame predictions so that these measurements can be repeated at stronger representations, where we expect the effect to be smaller and the analysis to be the same.
